% !TEX program = platex
% 数学ノート集・共通テンプレート v2（積分論2026年度第1部の版面を参考）
% A5判。章内で定義・定理・例・注意等の番号を共通化。
% コンパイル： platex -interaction=nonstopmode math_note_template_v2.tex (2回)
%              dvipdfmx math_note_template_v2.dvi
\documentclass[a5paper,10pt,onecolumn]{jsarticle}
\usepackage[dvipdfmx,top=22mm,bottom=20mm,inner=19mm,outer=18mm,headheight=13pt,headsep=6mm]{geometry}
\usepackage{amsmath,amssymb,amsthm,mathtools}
\usepackage{enumitem,fancyhdr,titlesec}
\usepackage[dvipdfmx,bookmarks=true,colorlinks=false]{hyperref}
\usepackage{pxjahyper}

% ----- 差し替える文書情報 -----
\newcommand{\NoteTitle}{補間と数値積分}
\newcommand{\NoteSeries}{数学ノート集}

% ----- 文字・組版 -----
% 参考PDFの質素な明朝＋Computer Modernの印象を引き継ぎながら，
% A5判で窮屈にならないよう本文の行間を少し確保する．
\linespread{1.13}
\setlength{\parindent}{1zw}
\setlength{\parskip}{0pt}
\setlength{\abovedisplayskip}{6pt plus 2pt minus 2pt}
\setlength{\belowdisplayskip}{6pt plus 2pt minus 2pt}
\setlength{\abovedisplayshortskip}{3pt plus 2pt minus 1pt}
\setlength{\belowdisplayshortskip}{4pt plus 2pt minus 1pt}
\setlist[enumerate]{topsep=3pt,itemsep=1pt,parsep=0pt,leftmargin=2.4zw}
\titleformat{\section}{\normalfont\normalsize\bfseries}{\thesection}{0.8zw}{}
\titlespacing*{\section}{0pt}{2.0ex plus .4ex}{.9ex}
\titleformat{\subsection}{\normalfont\normalsize\bfseries}{\thesubsection}{0.7zw}{}
\titlespacing*{\subsection}{0pt}{1.5ex plus .3ex}{.6ex}
\raggedbottom

% ----- 柱とページ番号：ページ上部、罫線なし -----
\fancypagestyle{mathnote}{%
  \fancyhf{}%
  \fancyhead[L]{\footnotesize \NoteSeries\ （\NoteTitle）}%
  \fancyhead[R]{\footnotesize \thepage}%
  \renewcommand{\headrulewidth}{0pt}%
  \renewcommand{\footrulewidth}{0pt}%
}
\pagestyle{mathnote}

% ----- 節内通し番号の定義・定理群 -----
\newtheoremstyle{mathnotetheorem}
 {4.5pt}{4.5pt}{\normalfont}{}%
 {\bfseries}{}{0.65zw}%
 {\thmname{#1}\thmnumber{ #2}\thmnote{（#3）}}
\theoremstyle{mathnotetheorem}
\newtheorem{theorem}{定理}[section]
\newtheorem{definition}[theorem]{定義}
\newtheorem{lemma}[theorem]{補題}
\newtheorem{proposition}[theorem]{命題}
\newtheorem{corollary}[theorem]{系}
\newtheorem{example}[theorem]{例}
\newtheorem{remark}[theorem]{注意}
\renewcommand{\proofname}{証明}
\renewcommand{\qedsymbol}{$\square$}
\makeatletter
\renewenvironment{proof}[1][\proofname]{%
  \par\pushQED{\qed}\normalfont\topsep4\p@\@plus2\p@\relax
  \trivlist\item[\hskip\labelsep\bfseries #1]\ignorespaces
}{\popQED\endtrivlist\@endpefalse}
\makeatother

% ----- 使うノートで適宜拡張する記号 -----
\newcommand{\R}{\mathbb{R}}
\newcommand{\norm}[1]{\lVert #1\rVert}
\newcommand{\abs}[1]{\lvert #1\rvert}

\begin{document}
\section*{到達目標と前提知識}
本稿では，指定された節点から補間多項式を構成し，
その誤差を証明に基づいて評価する．さらに
直交多項式からGauss求積の節点・重みを自分で計算し，
積分の近似値と誤差を検算できるようになることを目標とする．
前提は多項式の除法・微積分・Rolleの定理である．

\section{多項式補間}
$x_0,\ldots,x_n$ を相異なる実数とする．
\begin{theorem}[Lagrange補間]
任意の値 $y_0,\ldots,y_n$ に対し次数高々 $n$ の多項式 $p$ で
$p(x_i)=y_i$ を満たすものが一意に存在し，
\[
 p(x)=\sum_{i=0}^ny_i\ell_i(x),\qquad
 \ell_i(x)=\prod_{j\ne i}\frac{x-x_j}{x_i-x_j}
\]
で与えられる．
\end{theorem}
\begin{proof}
$\ell_i(x_j)=\delta_{ij}$ より表示式は指定値をとる．
二つの多項式の差は $n+1$ 個の相異なる根を持つ次数高々 $n$ の多項式なので零である．
\end{proof}
\begin{theorem}[補間誤差]
$f\in C^{n+1}([a,b])$ とし $x,x_i\in[a,b]$ とする．
補間多項式 $p$ についてある $\xi\in(a,b)$ が存在して
\[
 f(x)-p(x)=\frac{f^{(n+1)}(\xi)}{(n+1)!}
 \prod_{i=0}^n(x-x_i)
\]
となる．
\end{theorem}
\begin{proof}
$x=x_i$ なら明らかである．そうでない場合
$\omega(t)=\prod_i(t-x_i)$，$c=(f(x)-p(x))/\omega(x)$ とし
$g(t)=f(t)-p(t)-c\omega(t)$ と置く．
$g$ は $n+2$ 個の相異なる点で零になる．Rolleの定理を $n+1$ 回用いると
$g^{(n+1)}(\xi)=0$ となる．一方
$g^{(n+1)}=f^{(n+1)}-c(n+1)!$ なので式を得る．
\end{proof}
\section{Gauss求積の高精度性}
$w$ を $[a,b]$ 上ほとんど至る所正の可積分重みとする．
$\langle f,g\rangle=\int_a^b f(x)g(x)w(x)dx$ とする．
\begin{lemma}[直交多項式の零点]
$n$ 次のモニック多項式 $P_n$ が全ての次数 $n$ 未満の多項式と直交するとき，
$P_n$ は $(a,b)$ に $n$ 個の相異なる零点を持つ．
\end{lemma}
\begin{proof}
$P_n$ が $(a,b)$ に持つ奇数重複度の零点を $r$ 個とし，
それらを $c_1,\ldots,c_r$ とする．
$Q(x)=\prod_{j=1}^r(x-c_j)$ と置くと $P_nQ$ は
$(a,b)$ 上で符号が一定であり，零多項式ではない．
$w>0$ がほとんど至る所成立するから
$\int_a^bP_n(x)Q(x)w(x)dx\ne0$．
もし $r<n$ なら $\deg Q<n$ なので直交性に反する．
よって $r\ge n$．一方 $\deg P_n=n$ なので
$r=n$ であり，全零点は $(a,b)$ の相異なる単根である．
\end{proof}
\begin{theorem}[Gauss求積の代数的正確性]
モニック直交多項式 $P_n$ の $n$ 個の相異なる根
$x_1,\ldots,x_n\in(a,b)$ を用いて
$\lambda_i=\int_a^b\ell_i(x)w(x)dx$ と置く．
このとき
\[
 \int_a^bq(x)w(x)dx=\sum_{i=1}^n\lambda_iq(x_i)
\]
は次数高々 $2n-1$ の全多項式 $q$ に対して成立し，$\lambda_i>0$ である．
\end{theorem}
\begin{proof}
$q=P_ns+r$ と除算する（$\deg s,\deg r<n$）．
直交性から $\int P_nsw=0$，根 $x_i$ で $q(x_i)=r(x_i)$．
$r$ のLagrange補間式を積分して結論を得る．
さらに $\ell_i(x)^2-\ell_i(x)$ は各 $x_j$ で零になり，
次数高々 $2n-2$ だから上の正確性を用いて
$\lambda_i=\int\ell_i^2w>0$ である．
\end{proof}



\section{複合求積とGauss求積の誤差公式}
\begin{theorem}[複合台形公式の誤差]
$f\in C^2([a,b])$ とし，$N$ 等分点の幅を $h=(b-a)/N$ とする．
複合台形公式
\[
 T_N(f)=h\left(\frac{f(a)+f(b)}2+
 \sum_{j=1}^{N-1}f(a+jh)\right)
\]
について
\[
 \left|\int_a^bf(x)dx-T_N(f)\right|
 \le\frac{(b-a)h^2}{12}\norm{f''}_\infty
\]
が成立する．
\end{theorem}
\begin{proof}
各小区間 $[u,v]$ で $f$ の一次補間多項式を $p$ とすると
補間誤差公式から
$|f(x)-p(x)|\le\norm{f''}_\infty(x-u)(v-x)/2$．
これを積分すると
$|\int_u^v(f-p)|\le\norm{f''}_\infty h^3/12$．
$p$ の積分はその区間の台形公式に等しい．
全 $N$ 区間の誤差を加えると主張を得る．
\end{proof}
\begin{theorem}[Gauss求積の誤差公式]
$w$ を正の可積分重み，$P_n$ を前節のモニック直交多項式とする．
$f\in C^{2n}([a,b])$ に対しある $\xi\in(a,b)$ が存在して
\[
 \int_a^b f(x)w(x)dx-\sum_{i=1}^n\lambda_i f(x_i)
 =\frac{f^{(2n)}(\xi)}{(2n)!}
 \int_a^bP_n(x)^2w(x)dx
\]
となる．
\end{theorem}
\begin{proof}
各 $x_i$ で $f$ の値と一階導関数の値を一致させる
次数高々 $2n-1$ のHermite補間多項式 $H$ を取る．
存在と一意性は，その補間条件を全て零にする
次数 $2n-1$ 以下の多項式が $\prod_i(x-x_i)^2$ で
割り切れなければならず，従って零であることから従う．
重複節点での補間誤差公式（Rolleの定理を $2n$ 回適用して導く）から
各 $x$ に応じた $\xi_x\in(a,b)$ について
$ f(x)-H(x)=f^{(2n)}(\xi_x)P_n(x)^2/(2n)!$ である．
Gauss求積は $H$ に対して正確であり，$H(x_i)=f(x_i)$ だから
求積誤差は $\int(f-H)w$ に等しい．
$f^{(2n)}$ の最小値・最大値でこの積分を挟み，
$\int P_n^2w>0$ と中間値の定理を用いれば式を得る．
\end{proof}

\section{端点を指定するGauss--Radau求積}
\begin{theorem}[Gauss--Radau求積の存在と正値性]
$w$ を $[a,b]$ 上ほとんど至る所正の可積分関数とし，$n\ge2$ とする．
左端点 $a$ を必ず求積節点に含めるとき，
互いに異なる節点 $a=x_1<x_2<\cdots<x_n<b$ と
正の重み $\lambda_1,\ldots,\lambda_n$ が存在して
\[
 \int_a^b q(x)w(x)\,dx=\sum_{j=1}^n\lambda_j q(x_j)
\]
が次数 $2n-2$ 以下の全多項式 $q$ について成立する．
\end{theorem}
\begin{proof}
$(x-a)w(x)$ は $(a,b)$ 上ほとんど至る所正なので，
これを重みとする次数 $n-1$ のモニック直交多項式 $P_{n-1}$ が存在する．
存在は正定値な内積についてのGram--Schmidt直交化から従う．
前節の直交多項式の零点補題により，
$P_{n-1}$ の根 $x_2,\ldots,x_n$ は
$(a,b)$ 内の相異なる単根である．
節点を $x_1=a,x_2,\ldots,x_n$ とし，
次数 $n-1$ のLagrange基底を $\ell_j$ と置いて
$\lambda_j=\int_a^b\ell_j(x)w(x)\,dx$ と定義する．

次数 $2n-2$ 以下の $q$ を
\[
 q(x)=(x-a)P_{n-1}(x)s(x)+r(x),
 \quad\deg s\le n-2,\quad\deg r<n
\]
と除算する．直交性により
$\int_a^b(x-a)P_{n-1}(x)s(x)w(x)\,dx=0$ であり，
全節点で $(x_j-a)P_{n-1}(x_j)=0$ だから
$q(x_j)=r(x_j)$ である．
次数 $n-1$ 以下の $r$ は
$r=\sum_jr(x_j)\ell_j$ と補間されるので
求積式を得る．

最後に $\ell_j^2$ の次数は $2n-2$ 以下であり，
全節点で $\ell_j(x_i)^2=\ell_j(x_i)$ だから
\[
 \lambda_j=\sum_i\lambda_i\ell_j(x_i)^2
 =\int_a^b\ell_j(x)^2w(x)\,dx>0
\]
である．最後の不等号は $\ell_j$ が零でない多項式で
$w>0$ がほとんど至る所成り立つことから従う．
\end{proof}
\begin{remark}
Gauss--Radau公式は端点を一つ固定した $2n-2$ 次の代数的正確性を持つ．
$n$ 点のGauss公式に $n+1$ 点を加えるGauss--Kronrod拡張は
別問題であり，任意の正の重みについて正重みの拡張の存在を主張しない．
\end{remark}


\section{補間誤差と求積公式を計算する}
\begin{example}[補間誤差を厳密に算出する]\label{ex:interp-cubic}
$f(x)=x^3$ を節点 $-1,0,1$ で二次以下の多項式に
補間する．データ値は $-1,0,1$ である．
$f(x_i)=x_i$ が各節点で成り立つから，
Lagrange補間の一意性により $p(x)=x$ である．
従って実際の誤差は
\[
 f(x)-p(x)=x^3-x=(x+1)x(x-1).
\]
一方 $f^{(3)}(x)=6$ なので，一般の補間誤差公式は
\[
 \frac{f^{(3)}(\xi)}{3!}(x+1)x(x-1)
 =(x+1)x(x-1)
\]
を与え，完全に一致する．
一般に，補間節点での値が一致しても
区間の内部全体で正確とは限らないことが分かる．
\end{example}
\begin{proposition}[二点Gauss求積の節点と重み]\label{prop:gauss2}
重み $w(x)=1$ の区間 $[-1,1]$ に対して
二点Gauss求積公式は
\[
 \int_{-1}^1f(x)\,dx
 \ \simeq\
 f(-1/\sqrt3)+f(1/\sqrt3)
\]
である．この公式は次数3以下の多項式に対して正確である．
$f(x)=x^4$ の場合の誤差は厳密に $8/45$ となる．
\end{proposition}
\begin{proof}
二次のモニック直交多項式は
$P_2(x)=x^2-1/3$ である．
実際，$\int_{-1}^1P_2(x)\,dx=0$ であり，
$\int_{-1}^1xP_2(x)\,dx=0$ は奇関数性による．
零点は $\pm1/\sqrt3$ である．
二点におけるLagrange補間基底 $\ell_\pm$ の積分は
対称性により等しく，その和は
$\int_{-1}^1 1\,dx=2$ だから各重みは1である．
従って先のGauss求積定理により3次以下で正確である．
$x^4$ の積分は $2/5$，求積値は
$2(1/\sqrt3)^4=2/9$ なので差は $8/45$．
さらに一般誤差公式では
$f^{(4)}=24$，
$\int_{-1}^1P_2^2\,dx=8/45$ であり
同じ誤差を与える．
\end{proof}
\section*{到達確認問題と完全解答}
\noindent\textbf{問題．}
$f(x)=x^2$ を $[0,1]$ で積分するとき，
等間隔の二分割による複合台形公式の値と厳密積分値を求め，
誤差を評価定理と比較せよ．

\medskip\noindent\textbf{解答．}
$h=1/2$ で
\[
 T_2(f)=\frac12\left(
 \frac{f(0)+f(1)}2+f(1/2)\right)
 =\frac12\left(\frac12+\frac14\right)=\frac38.
\]
厳密積分値は $\int_0^1x^2dx=1/3$ であり，
絶対誤差は
\[
 \frac38-\frac13=\frac1{24}.
\]
$f''=2$ だから一般の誤差上界は
$(b-a)h^2\|f''\|_\infty/12
 =1\cdot(1/4)\cdot2/12=1/24$．
従ってこの例では上界に等号が成立する．

\end{document}
